%%%%%%%%%%%%%%%%%%%%%%%%%%%%%%%%%%%%%%%%%
% Important note:
% This template requires the resume.cls file to be in the same directory as the
% .tex file. The resume.cls file provides the resume style used for structuring the
% document.
%
%%%%%%%%%%%%%%%%%%%%%%%%%%%%%%%%%%%%%%%%%

%----------------------------------------------------------------------------------------
%	PACKAGES AND OTHER DOCUMENT CONFIGURATIONS
%----------------------------------------------------------------------------------------

\documentclass{resume} % Use the custom resume.cls style
\usepackage{setspace}
\usepackage{microtype}
\usepackage{enumitem}
\usepackage{bbding}
\usepackage{microtype}
\usepackage{latexsym}
\usepackage{amsfonts}

\usepackage[left=0.75in,top=0.6in,right=0.75in,bottom=0.6in]{geometry} % Document margins
\newcommand{\tab}[1]{\hspace{.2667\textwidth}\rlap{#1}}
\newcommand{\itab}[1]{\hspace{0em}\rlap{#1}}

\newlist{myinneritemize}{itemize}{1}
\setlist[myinneritemize]{label=$\bullet$, itemsep=2pt, topsep=2pt, partopsep=0pt}


\name{Zeyi Liao} % Your name
\vspace{-10pt}
\address{\small {https://lzy37ld.github.io}} % Your address
%\address{123 Pleasant Lane \\ City, State 12345} % Your secondary addess (optional)
\address{Phone:(+1) 213 551 8327\\ Email:lzy37ld@gmail.com} % Your phone number and email

\begin{document}

%----------------------------------------------------------------------------------------
%	EDUCATION SECTION
%----------------------------------------------------------------------------------------

\begin{rSection}{Education}
%--copy and paste this region  if you need more--
\setstretch{0.6} % 设置行间距为1.0倍

{\bf Ohio State University} \hfill {\em 2023 - 2028 (expected)} 
\\ \textit{Ph.D in Computer Science} \hfill {Advisor: Prof.Huan Sun}
\\

{\bf University of Southern California} \hfill {\em 2021 - 2023 } 
\\ \textit{Master in Electronic and Computer Engineering} \hfill {Advisor: Prof.Xiang Ren}
\\

{\bf Beijing Jiaotong Univeristy} \hfill {\em 2016 - 2020 } 
\\ \textit{Bachelor in Electronic Engineering} \hfill {Advisor: Prof.Yishuai Chen}
\\
%--copy and paste this region  if you need more--
\end{rSection}

%----------------------------------------------------------------------------------------
%	EXPERIENCE SECTION
%----------------------------------------------------------------------------------------
\begin{rSection}{Publications}
%--copy and paste this region  if you need more--
% {\bf RobustLR: A diagnostic benchmark for evaluating logical robustness of deductive reasoners}\\
% description of the experience\\\\

\leftskip=0pt % 负值用于将左边界向左移动，可以根据需要微调这个值
\setstretch{0.6}

- {\bf \textls[-50]{In Search of the Long-tail: Systematic Generation of Long-tail Knowledge via Logical Rule Guided Search}}
\vspace{-1pt}
\begin{myinneritemize}
    \item Huihan Li, \underline{Zeyi Liao}, Yuting ning, Siyuan Wang, Xiang Lorraine Li, Ximing Lu, Faze Brahman, Yejing Choi, Xiang Ren
    \vspace{-3pt}
    \item \textit{Under Review}
    \vspace{-3pt}
    \item We propose a pipeline to efficiently generate long-tail knowledge while being factually correct and reliable.
\end{myinneritemize}


- {\bf \textls[-50]{ChatCounselor: A Large Language Models for Mental Health Support}}
\vspace{-3pt}
\begin{myinneritemize}
    \item June M Liu, Donghao Li, He Cao, Tianhe Ren, \underline{Zeyi Liao}, Jiamin Wu
    \vspace{-3pt}
    \item \textit{CIKM 2023}
    \vspace{-3pt}
    \item We present ChatCounselor, a large language model (LLM)
solution designed to provide mental health support.
\end{myinneritemize}

- {\bf \textls[-50]{RobustLR: A diagnostic benchmark for evaluating logical robustness of deductive reasoners}}
\vspace{-3pt}
\begin{myinneritemize}
    \item Soumya Sanyal, \underline{Zeyi Liao}, Xiang Ren
    \vspace{-3pt}
    \item \textit{EMNLP 2022}
    \vspace{-3pt}
    \item We propose a benchmark to diagnose the deductive reasoning of LMs.
\end{myinneritemize}


    

%--copy and paste this region  if you need more--
\end{rSection}

\begin{rSection}{Ongoing Projects}

- \textbf{\textls[-50]{IRL gains extra supervision compared to RLHF}}\\
\textit{Desc}: We aim to demonstrate that IRL can benefit from additional forms of supervision compared to RLHF. To illustrate, we collect preference data from HH and categorize pairs with significant differences as the "demonstration, trajectory" group, while the rest fall into the regular preference group. Our preliminary findings indicate that, under the same data constraints, IRL and IRL + RLHF yield superior results.\\
\vspace{-10pt} \\
- \textbf{\textls[-50]{More robust Attribution Evaluation}}\\
\textit{Desc}: Although Language Models (LLMs) are susceptible to generating hallucinations, quantifying this issue automatically is challenging, and assessing it through human evaluation is time-consuming and costly. Despite previous efforts to enhance machine evaluation for this problem, they still fall short of human performance. Our goal is to introduce a new evaluator tuned on a diverse curated dataset and propose utilizing the entire web content as input, rather than solely relying on human-annotated data, which we frequently find to be uninformative.



\end{rSection}
\begin{rSection}{Awards and Scholarships} 
{\bf Master Honor Student}{, University of Souther California
} \hfill{\em 08/22}\\
{\bf The Third Prize in Toxic classification}{, Southern Grid Company
} \hfill{\em 12/2020}\\
{\bf First-class University Scholarship}{, Beijing Jiaotong University
} \hfill{\em 11/2020,2019}\\
{\bf Honorable Mention}{, Interdisciplinary Contest in Modeling Certificate of Achievement
} \hfill{\em 04/2019}\\
{\bf Merit Student}{, Beijing Jiaotong University
} \hfill{\em 12/2017}\\
\end{rSection}


\end{document}----------------------------

